\begin{abstract}
\noindent
Layered queueing networks (LQNs) extend classical queueing networks by 
modelling synchronous calls between layered services and the blocking they induce, enabling the analysis of distributed software systems. 
LQNS and LINE's \texttt{SolverLN} are two general-purpose tools for solving LQN models, however they have become bloated over years of development.
This
report describes the construction of \texttt{SolverLNSimple},
a simplified LQN solver, built as the minimal solution under the paradigm of only adding necessary features based on the set of models considered. The build proceeds through a series of features, each triggered by a specific case of LQN model the existing design failed to handle: ensemble prefix coupling, graph-derived coupling, multi-class models, an MVA-kernel overhaul, larger fan-out callers and activity graph integrations. 
This bottom-up discipline restricts the machinery added, providing a lightweight, efficient solver for a large class of LQN models. Once the solver was sufficiently correct, its non-generality was exploited to discover performance benefits.
An evaluation of the solver, on a 78-fixture corpus, against \texttt{SolverLN} and LQNS verifies its correctness and characterises its performance, where we see a large performance win whilst maintaining accuracy on the set of models the solver claims to handle. 
73.6\% of the iteration count difference is attributable to \texttt{SolverLN}'s \texttt{iter\_min} convergence floor, with the rest reflecting genuine algorithmic difference, however the omission of a moving-average convergence detection and interlock machinery carries a small, quantified accuracy cost.
Regarding runtime, much of the difference is attributed to \texttt{SolverLN}'s interlock correction, where cost grows with fan-out breadth, and a computed-but-not-consumed task-caller traversal. 
\end{abstract}
